\clearpage
\thispagestyle{empty}

\begin{center}
	\vspace*{4cm}
	{\Huge \textbf{KẾT THÚC PHẦN 1}} \\[1.5cm]
	
	\rule{0.8\textwidth}{0.4pt} \\[0.8cm]
	
	\parbox{0.9\textwidth}{
		\centering
		\textit{
			Phần 1 của giáo trình đã cung cấp cho bạn một nền tảng lý thuyết toàn diện, 
			từ những nguyên lý ngôn ngữ học cơ bản, các mô hình thống kê, 
			các kiến trúc nơ-ron kinh điển, cho đến sự thống trị của Transformer 
			và các kỹ thuật tiên tiến để tinh chỉnh, căn chỉnh, và đánh giá các 
			Mô hình Ngôn ngữ Lớn. Với nền tảng này, bạn đã sẵn sàng để bước sang 
			Phần 2, nơi chúng ta sẽ biến những lý thuyết này thành các kỹ năng và 
			quy trình thực chiến để xây dựng các ứng dụng NLP thực tế.
		}
	} \\[1cm]
	
	\rule{0.8\textwidth}{0.4pt} \\[2cm]
\end{center}

\noindent
\textbf{Tóm lược nội dung Phần 1:}
\begin{itemize}
	\item \textbf{Chương 1 – Nhập môn và Nguyên lý nền tảng:} 
	Định nghĩa NLP, mối quan hệ với AI, các kỷ nguyên phát triển, 
	kiến thức ngôn ngữ học, toán học, và các vấn đề đạo đức trong NLP.
	\item \textbf{Chương 2 – Biểu diễn văn bản và Mô hình thống kê:} 
	Các kỹ thuật BoW, TF-IDF, N-gram, smoothing, 
	mô hình phân loại, học chuỗi, và LDA.
	\item \textbf{Chương 3 – Các kiến trúc mạng nơ-ron kinh điển:} 
	Word embeddings (Word2Vec, GloVe, FastText), 
	autoencoder, RNN/LSTM/GRU, seq2seq với attention, 
	CNN cho NLP, contextualized embeddings, và các mô hình sinh pre-Transformer.
	\item \textbf{Chương 4 – Kỷ nguyên Transformer và LLMs:}
	Kiến trúc Transformer gốc, phân loại LLMs, tokenization hiện đại,
	ngữ cảnh dài (FlashAttention, YaRN, Ring Attention), MoE (Switch, DeepSeekMoE),
	khối decoder hiện đại (GQA, MLA, MTP).
	\item \textbf{Chương 5 – Tiền huấn luyện LLM:}
	Scaling laws (Kaplan, Chinchilla), dữ liệu tiền huấn luyện (The Pile, RefinedWeb, FineWeb, DCLM, DoReMi, dữ liệu tổng hợp),
	bộ tối ưu (AdamW, Sophia, Adam-mini, Muon), công thức Llama 3 và DeepSeek-V3.
	\item \textbf{Chương 6 – Tinh chỉnh và Căn chỉnh mô hình:}
	Fine-tuning và thích ứng miền, PEFT (LoRA, DoRA, QLoRA, Adapter, IA$^3$), prompt engineering,
	instruction tuning, RLHF (InstructGPT), DPO, KTO, ORPO, GRPO.
	\item \textbf{Chương 7 – Mô hình suy luận và tính toán lúc suy luận:}
	Verifier (ORM, PRM), scaling test-time compute, STaR/ReST$^{EM}$/Quiet-STaR,
	học tăng cường với phần thưởng kiểm chứng được (DeepSeek-R1, DAPO, GSPO).
	\item \textbf{Chương 8 – Tìm kiếm và Sinh Tăng cường (RAG):}
	Pipeline RAG từ cơ bản đến nâng cao, biểu diễn chuyên biệt cho truy xuất, RAG đa bước và RAG dạng tác tử.
	\item \textbf{Chương 9 – Hiệu quả LLM:}
	Chưng cất (MiniLLM, GKD, Distilling Step-by-Step, KARD, RLAD), lượng tử hóa (LLM.int8, SmoothQuant, GPTQ, AWQ),
	PagedAttention, giải mã suy đoán, phục vụ tách rời prefill/decode.
	\item \textbf{Chương 10 – Hệ thống và Kiến trúc nâng cao:}
	Mô hình đa phương thức (CLIP, LLaVA, Gemini), hệ thống tác tử, world models.
	\item \textbf{Chương 11 – Đánh giá và Benchmark:}
	Các metric kinh điển, benchmark (GLUE, MMLU, GSM8K, MATH, TruthfulQA, HELM),
	thách thức trong đánh giá, LLM-as-a-Judge và Chatbot Arena,
	đánh giá tác tử (AgentBench, GAIA, SWE-bench) và tính tái lập.
\end{itemize}